\documentclass{article}
\title{Advances in Attention Mechanisms}
\author{Jane Doe}
\begin{document}
\maketitle

\section{Introduction}
Attention mechanisms have revolutionised sequence modelling by allowing models
to focus selectively on relevant parts of the input. This work examines recent
advances in multi-head attention and sparse attention variants.

\section{Background}
The original attention mechanism was introduced in the context of neural machine
translation. \textbf{Self-attention} extends this idea by computing attention
weights within a single sequence rather than between an encoder and decoder.

\subsection{Multi-Head Attention}
Multi-head attention projects the input into multiple subspaces, enabling the
model to capture diverse relational patterns. Each head operates independently
and the outputs are concatenated before a final linear projection.

\section{Experiments}
We benchmarked three attention variants on the \textit{WMT-14} translation
dataset. Sparse attention reduced memory consumption by 40\% with less than
1\% degradation in BLEU score.

\section{Conclusion}
Efficient attention variants make it practical to apply transformer models
to longer sequences without prohibitive memory costs.

\end{document}
